\section{Choosing what to reuse from Causal-Copilot}
\label{sec:exec25-copilot-reuse}

A bank evaluating Causal-Copilot needs a component decision as well as the
scientific assessment in Section~\ref{exec17:copilot-development}. The useful
question is which responsibilities the proposed library should retain and
which it could delegate after a bounded trial. The inspected README links the
paper, repository and demonstration under one project identity and describes
V1 and V2 releases. Those names establish a documented project lineage; they
do not identify three independently validated engines. This assessment binds
the Lancelot39 repository at commit \texttt{a31aa7c}, dated 15 May 2026, and the
21 April 2025 paper separately. No claim that the later source produced the
earlier experiments follows from their shared name. Other similarly named
repositories require their own identity and revision checks.
\citep{exec17_copilot_paper,exec17_copilot_code}

The present choice is architectural reference use. The workflow decomposition
is a concrete comparator for business dialogue, method selection, execution
and reporting. Direct adoption would additionally require a reproducible
environment, a complete dependency and rights inventory, and demonstrated
semantic behavior. Selective adaptation of a parser or report template may be
less extensive, but still requires checking what information it loses. An
isolated integration could retain the upstream workflow while placing an
independent contract around it. A language port would assume responsibility
for every numerical, graph and failure conversion; it would not by itself
repair causal authority or reduce maintenance obligations. The proposed
library should reject automatic promotion of a suggested direction to an
identified intervention claim, while retaining useful components whose
contracts can be qualified. These are component choices, not a declaration
that the entire project is either safe or unsuitable.
\citep{exec17_copilot_code,exec25_copilot_choice}

\subsection{The boundary between assistance and scientific authority}

Consider again a retention offer assigned at a stated deadline. A proposed
request would carry a stable request identifier, a data and schema digest,
the eligible population, assignment and outcome clocks, treatment versions,
the requested contrast and a ledger of assumptions with their provenance.
The agent may suggest revisions, but a revision changes a versioned request
rather than silently changing the estimand. The selected upstream filter
bypasses algorithm filtering when the user already selected an algorithm;
that preference therefore needs a separate assumption check at the proposed
library boundary. A successful method call must not certify the assumptions
that motivated it.~\citep{exec17_copilot_code}

The corresponding result would distinguish an observational graph, a causal
claim conditional on assumptions, and an unavailable answer. It would preserve
variable identities, the convention for matrix entries, endpoint marks, lags
and raw versus revised graphs. The selected evaluator uses the convention
that entry $(i,j)$ refers to the arrow $j\to i$; silently adopting the
opposite convention would reverse a reported intervention direction. The
earlier source review also found that the inspected judge calls the bootstrap
helper with \texttt{ts=False}. A declared temporal request therefore needs
evidence that its actual path preserves the intended dependence and lag
semantics before receiving a temporal uncertainty claim. Neither a stored
lag array nor a successful JSON response supplies that evidence.
\citep{exec17_copilot_code}

For a minimal proposed representation, let $\mathcal G(r)$ denote the set of
directed possibilities represented by result $r$. With just two variables,
an unresolved adjacency can represent
\[
 \mathcal G(r_u)=\{X\to Y,Y\to X\},\qquad
 \mathcal G(r_d)=\{X\to Y\},\qquad
 \mathcal G(r_0)=\{\hbox{no edge}\}.
\]
The first two have the same skeleton but different directional information;
the first and third distinguish unknown direction from an absent relation.
This small authored representation is not a full CPDAG or PAG definition. It
demonstrates why a conversion test should compare represented possibilities,
not merely dimensions or adjacency counts. The complete adapter would need
the appropriate graph-class semantics, including latent-confounding marks
where relevant.~\citep{exec25_copilot_choice}

An edit proposal should carry the preceding graph digest, a reason, its
evidence, its author and the exact endpoints it changes. Scientific authority
and runtime permission are separate. A reviewer may permit an edit without
certifying the new direction. Conversely, a well-motivated edit must still
fail if it targets an outdated graph. The inspected refinement path and its
conditional witness motivate a symmetric protected-pair rule: both orderings
of a protected pair must be checked. An unknown answer should preserve
uncertainty rather than become a no-edge decision. The protected-pair rule
implements a declared authority policy; it does not prove that the protected
edge is true. Later statistical pruning needs its own explicit precedence
rule and audit, rather than an implicit exemption.
\citep{exec17_copilot_code,exec17_copilot_exact,exec25_copilot_choice}

\subsection{A public trial with distinct scientific and operational oracles}

The next useful experiment is a bounded prototype on generated mechanisms
and permitted public data. Before execution, the study should freeze task
families, graph conventions, provider and package versions, prompts, model
settings, seeds, splits and scoring rules. Generated structural models supply
known intervention effects for scientific checks. Public observational data
can test missingness, schema changes and workflow behavior, but unknown
causal truth must remain unknown. The earlier observational-equivalence
construction supplies a refusal test: without additional distinguishing
information, the two admitted worlds cannot justify one unconditional effect.
\citep{exec17_copilot_exact,exec25_copilot_choice}

Three comparisons answer different questions. A fixed method with a
prespecified configuration tests the cost and benefit of adaptive selection.
The source-pinned upstream workflow, once permitted and reproducible, tests
the documented agent journey. A proposed contract-wrapped workflow tests
whether preserved ambiguity, edit authority and typed refusals improve the
decision under the same tasks and information. A human-assisted baseline
can measure the review burden, with effort and expertise recorded. These
arms must receive matched data and information; allowing only one arm to use
the true graph would answer a different question. The inspected evaluator's
truth-assisted orientation may be retained as a named compatibility score,
but it cannot be the sole measure of an operational direction prediction.
\citep{exec17_copilot_code,exec25_copilot_choice}

The proposed qualification suite separates several failure modes. A variable
permutation fixture checks names and matrix conventions. Opposite admissible
orientations test ambiguity preservation. Paired lag-one and lag-five
requests test whether the result retains the information clock. Protected
edits in both identifier orders and a stale graph digest test edit authority.
Malformed answers, provider refusal, timeout, cancellation, duplicate delivery
and a crash between execution and report generation test the runtime
contract. Repeating a request after an uncertain timeout must not invent a
completion or double-count an accepted result. These are proposed tests; no
upstream package or provider was run to pass them in this assessment.
\citep{exec25_copilot_choice}

Scientific acceptance should be prespecified separately from those exact
interface invariants. A structurally valid graph can still be wrong, and
well-calibrated uncertainty in one stipulated regime need not transfer to
another. The trial should report graph ambiguity, target identification,
effect error and coverage where ground truth is available, together with
correct refusals, inappropriate refusals and failed runs. Concrete error
tolerances, power and sample size require the selected decision and precision
requirements; they are not supplied here by an arbitrary universal threshold.
This design can establish whether a larger experiment is worth its cost
without claiming authentic bank validation.~\citep{exec25_copilot_choice}

\subsection{Rights, ownership and the denominator of value}

The inspected root and vendored causal-learn notices are MIT notices. They
support reference and component assessment subject to their conditions, but
do not replace an inventory of all dependencies, model services, assets and
data. The selected requirements mix pinned packages with an unpinned
\texttt{huggingface-hub}; default pytest settings exclude integration, LLM
and LaTeX markers. Thus a named historical commit is not yet a complete
reproducible or maintained deployment profile. Present upstream support and
compatibility are unresolved rather than inferred from a historical README.
\citep{exec17_copilot_code,exec17_copilot_rights}

The scientific owner would sign the target, assumptions, graph class and
scoring rule. The adapter owner would maintain conversion, edit and failure
contracts. Platform and procurement owners would establish the environment,
service terms and retention policy. An evaluation owner would maintain
independent oracles and all-request accounting. A port additionally needs a
named maintainer for semantic parity across upstream changes. These are
proposed responsibilities; no institution has accepted them merely because
they appear in this proposal.~\citep{exec25_copilot_choice}

For an illustrative trial, suppose preparation costs USD\,1,200, all attempts
USD\,1,800, review USD\,2,400 and allocated maintenance USD\,600. Total cost
is USD\,6,000. With 120 requests, 100 completions and 80 scientifically
admissible results, the costs per request, completion and admissible result
are respectively USD\,50, USD\,60 and USD\,75. Completion is $5/6$, whereas
the admissible fraction is $2/3$. These exact authored quantities illustrate
why reporting only successful calls hides the cost of failures and refusals.
They are neither a vendor quote nor an estimated return on bank investment.
The useful denominator must follow the actual decision, with every attempt
retained and admissibility criteria fixed before outcomes are observed.
\citep{exec25_copilot_choice}

Reference use can therefore be decided now without declaring integration
ready. Reconsider adoption only after a selected artifact and complete
environment have passed the stated scientific and runtime contracts, rights
are established for the intended deployment, and matched full-cost evidence
supports the decision. Experiment-to-code binding remains necessary for any
reproduction claim. Generic agent repositories and diagram assistants still
require their own concrete candidate decisions.
\label{sec:exec25-copilot-reuse-end}

